\documentclass[letter, 10pt]{article}

\usepackage{amsmath}
\usepackage{xcolor}
\usepackage{subfig}
\usepackage[margin = 1in]{geometry}
\usepackage{graphicx}
\usepackage{epstopdf}
\usepackage{pxfonts}
\usepackage{setspace}
\usepackage{titlesec}
\usepackage{booktabs}
\usepackage{fancyhdr}
\usepackage{listings}
\usepackage{multirow}
\usepackage{tikz}
\usepackage{ragged2e}

\usetikzlibrary{positioning}
\usetikzlibrary{fit}

\definecolor{evolotus-color}{RGB}{56,64,96}

\titlespacing{\section}{0pt}{\baselineskip}{0.5\baselineskip}
\titleformat{\section}{\normalfont\fontsize{14}{0}\bfseries}{\thesection}{1em}{}
\titlespacing{\subsection}{0pt}{\baselineskip}{0.2\baselineskip}
\titleformat{\subsection}{\normalfont\fontsize{12}{0}\bfseries}{\thesubsection}{1em}{}

\def\helv{\fontfamily{phv}\bfseries\selectfont}

\begin{document}
\begin{onehalfspacing}

\begin{titlepage}

	\noindent\includegraphics[scale = 0.4]{evolutus-with-text-blue.eps}
	\vspace{25\baselineskip}\vfill
	
	{\noindent\Large\uppercase{16-662 Robot Autonomy}\\[6pt]
	\Huge\uppercase{\bf Project 2A: Report}}\\[6pt]
	
	\noindent Team Evolutus \\
	\noindent William Ku, Yuhan Long, Yu-Te Cheng, Dawei Wang. \vspace{12pt}
	
	\noindent\today\vfill

\end{titlepage}

\pagestyle{fancy}
\lhead{Project 2B}
\rhead{Team Evolutus}

\setlength{\parskip}{0.5\baselineskip}
\RaggedRight
\parindent=2em

\newpage

\section{Unscented Kalman Filter}

Here we presented fours sets of the our Unscented Kalman Filter results. Figure \ref{baseline} shows the result of our baseline test. In the baseline test, we set all the noise covariances number as 0.001. To demonstrate the effects of changing the initial noise sigma value, IMU noise value and observation noise value, we change to each group of sigma value to make them much larger than the baseline. 
\begin{figure}[h]
	\centering
	\begin{minipage}{.25\textwidth} 
  		\centering
    		\includegraphics[width=1\linewidth]{b_x.eps}
	\end{minipage}%	
	\begin{minipage}{.25\textwidth} 
		\centering
    		\includegraphics[width=1\linewidth]{b_y.eps}
	\end{minipage}%
	\begin{minipage}{.25\textwidth}  		
		\centering
    		\includegraphics[width=1\linewidth]{b_z.eps}
	\end{minipage}%
	\begin{minipage}{.25\textwidth}  
		\centering
    		\includegraphics[width=1\linewidth]{b_yaw.eps}
	\end{minipage}
	\centering
	\caption{Baseline results with all parameters set to 0.001.}
	\label{baseline}
\end{figure}

\begin{figure}[h]
	\centering
	\begin{minipage}{.25\textwidth} 
  		\centering
    		\includegraphics[width=1\linewidth]{i_x.eps}
	\end{minipage}%	
	\begin{minipage}{.25\textwidth} 
		\centering
    		\includegraphics[width=1\linewidth]{i_y.eps}
	\end{minipage}%
	\begin{minipage}{.25\textwidth}  		
		\centering
    		\includegraphics[width=1\linewidth]{i_z.eps}
	\end{minipage}%
	\begin{minipage}{.25\textwidth}  
		\centering
    		\includegraphics[width=1\linewidth]{i_yaw.eps}
	\end{minipage}
	\centering
	\caption{Test results with 100 times initial noise coefficients.}
	\label{highi}
\end{figure}

\begin{figure}[h]
	\centering
	\begin{minipage}{.25\textwidth} 
  		\centering
    		\includegraphics[width=1\linewidth]{q_x.eps}
	\end{minipage}%	
	\begin{minipage}{.25\textwidth} 
		\centering
    		\includegraphics[width=1\linewidth]{q_y.eps}
	\end{minipage}%
	\begin{minipage}{.25\textwidth}  		
		\centering
    		\includegraphics[width=1\linewidth]{q_z.eps}
	\end{minipage}%
	\begin{minipage}{.25\textwidth}  
		\centering
    		\includegraphics[width=1\linewidth]{q_yaw.eps}
	\end{minipage}
	\centering
	\caption{Test results with 50 times IMU noise coefficients.}
	\label{highq}
\end{figure}

\begin{figure}[h]
	\centering
	\begin{minipage}{.25\textwidth} 
  		\centering
    		\includegraphics[width=1\linewidth]{r_x.eps}
	\end{minipage}%	
	\begin{minipage}{.25\textwidth} 
		\centering
    		\includegraphics[width=1\linewidth]{r_y.eps}
	\end{minipage}%
	\begin{minipage}{.25\textwidth}  		
		\centering
    		\includegraphics[width=1\linewidth]{r_z.eps}
	\end{minipage}%
	\begin{minipage}{.25\textwidth}  
		\centering
    		\includegraphics[width=1\linewidth]{r_yaw.eps}
	\end{minipage}
	\centering
	\caption{Test results with 100 times observation noise coefficients.}
	\label{highr}
\end{figure}

Figure \ref{highi} shows the UKF results when we increase the initial noise sigma value by 100 times. From the result, we know that when the initial noise value is increased, the UKF will have a large error initially, but will return tracking after several second. For some fast changing state (z position in this case), the UKF needs longer time to reduce the estimation errors. 

Figure \ref{highq} shows the UKF results when we increase the IMU noise value by 50 times. Changing the IMU noise and IMU bias noise will increase the UKF error for in the model update process. In the correction process, the error will be corrected by the observations. Therefore, adding IMU noise will make the UKF output noisy, but will not create a large bias of the results.

Figure \ref{highr} shows the UKF results when we increase the observation noise by 100 times. Increasing the observation noise will not create noise on the UKF output, but it can create bias on the result and affect the accuracy. This is because a wrong observation will lead to a bias correction that will affect all the state vector. 

We analysis the covariance of the each variable in the process model and the observation model to estimate our noise parameters. The observation noise is estimated by taking the mean values from the ground truth data and calculating the covariance of the observation results. The IMU bias noise is provided in the class as 0.03. The IMU acceleration noise and angular velocity noise is estimated by calculating covariances within a relatively static time of the drone. We tuned the initial sigma value manually. It shows that as long as the initial values are kept within a reasonable range, it will not affect the result too much. The parameters are shown in Table \ref{np}. The UKF output after noise parameter tuning is shown in Figure \ref{tuning}.

\begin{center}
\begin{table}[h]
\caption{Noise Parameters}
\label{np}
\centering\begin{tabular}{cccc}
\hline
init\_pos\_sigma    & 0.01    & sigma\_bw           & 0.03      \\
init\_rpy\_sigma    & 0.001   &  sigma\_a            & 0.08      \\
init\_vel\_sigma    & 0.001    &  sigma\_ba           & 0.03    \\
init\_bw\_sigma     & 0.001    &  sigma\_xy           & 0.077     \\
init\_ba\_sigma     & 0.001      &  sigma\_z            & 0.05  \\
sigma\_w            & 0.03       & sigma\_ yaw           & 0.02  \\ \hline
\end{tabular}
\end{table}
\end{center}

\begin{figure}[h]
	\centering
	\begin{minipage}{.25\textwidth} 
  		\centering
    		\includegraphics[width=1\linewidth]{t_x.eps}
	\end{minipage}%	
	\begin{minipage}{.25\textwidth} 
		\centering
    		\includegraphics[width=1\linewidth]{t_y.eps}
	\end{minipage}%
	\begin{minipage}{.25\textwidth}  		
		\centering
    		\includegraphics[width=1\linewidth]{t_z.eps}
	\end{minipage}%
	\begin{minipage}{.25\textwidth}  
		\centering
    		\includegraphics[width=1\linewidth]{t_yaw.eps}
	\end{minipage}
	\centering
	\caption{Test results after noise value tuning.}
	\label{tuning}
\end{figure}

We calculate the mean square error of the UKF estimate position against the ground truth position. Figure \ref{MSE} show the error plot and the error distribution. From the distribution, we can see that after tuning the UKF noise parameter, the largest Mean Square Error is 0.013. The mean error value is 0.0008. Most of the error value is below 0.004. 

\begin{figure}[h]
	\centering
	\begin{minipage}{.5\textwidth} 
  		\centering
    		\includegraphics[width=0.8\linewidth]{MSE.eps}
	\end{minipage}%	
	\begin{minipage}{.5\textwidth} 
		\centering
    		\includegraphics[width=0.8\linewidth]{MSE_dist.eps}
	\end{minipage}
	\centering
	\caption{Mean Square Error of UKF position estimate plot and corresponding distribution.}
	\label{MSE}	
\end{figure}

We select the time interval from 15s to 20s to compare the empirical standard deviation and estimate UKF 1-sigma value. Figure \ref{MSEinterval} shows position error vs. time plot and error distribution on x, y, z direction. The empirical standard deviation in this time interval is: std\_x: 0.026569 , std\_y: 0.015594, std\_z: 0.016420. The estimated UKF 1-sigma values in this time is x: 0.030517 ,y: 0.030510, z: 0.014851. The empirical standard deviation is smaller than the UKF estimation result. 

\begin{figure}[h]
	\centering
	\begin{minipage}{.5\textwidth} 
  		\centering
    		\includegraphics[width=0.8\linewidth]{MSE_t.eps}
	\end{minipage}%	
	\begin{minipage}{.5\textwidth} 
		\centering
    		\includegraphics[width=0.8\linewidth]{MSE_x.eps}
	\end{minipage}

	\centering
	\begin{minipage}{.5\textwidth} 
  		\centering
    		\includegraphics[width=0.8\linewidth]{MSE_y.eps}
	\end{minipage}%	
	\begin{minipage}{.5\textwidth} 
		\centering
    		\includegraphics[width=0.8\linewidth]{MSE_z.eps}
	\end{minipage}
	\centering
	\caption{Position Error vs. Time and each error distribution.}
	\label{MSEinterval}	
\end{figure}

\vfill
\begin{center}
	\includegraphics[scale = 0.4]{evolutus-black.eps}
\end{center}


\thispagestyle{fancy}
	
\end{onehalfspacing}
\end{document} 